\documentclass[a4paper, serif, 10pt]{moderncv}

%% moderncv
% color and style
\moderncvcolor{black}
\moderncvstyle{banking}

%% page layout/margins
\usepackage[top=2.5cm, bottom=2.5cm, left=1.5cm, right=1.5cm]{geometry}

\nopagenumbers{}

%% Babel config for French language
% ref:
% - https://latex3.github.io/babel/guides/locale-french.html
\usepackage[french]{babel}

%% fontspec
\usepackage{fontspec}

\IfFontExistsTF{Linux Libertine}{
  \setmainfont[Ligatures={TeX, Common}, Numbers=OldStyle]{Linux Libertine}
}{}
\IfFontExistsTF{Linux Libertine O}{
  \setmainfont[Ligatures={TeX, Common}, Numbers=OldStyle]{Linux Libertine O}
}{}

%% CTeX
% CJK support, used to show CJK characters in the resume
%
% - heading=false: disable ctex heading
% - punct=kaiming: use kaiming punctuations styles for CJK
% - scheme=plain: use plain scheme, do not override \normalsize font size
% - space=auto: space settings for CJK characters
%
% CJK fonts are resolved per locale language, e.g. Noto Serif/Sans CJK JP
% for Japanese resumes, with fallback to the next region if not installed
%
% ref:
% - http://ctan.mirrorcatalogs.com/language/chinese/ctex/ctex.pdf
\usepackage[UTF8, heading=false, punct=kaiming, scheme=plain, space=auto]{ctex}

\IfFontExistsTF{Noto Serif CJK SC}{
  \setCJKmainfont{Noto Serif CJK SC}
}{}
\IfFontExistsTF{Noto Sans CJK SC}{
  \setCJKsansfont{Noto Sans CJK SC}
}{}

%% line spacing
\usepackage{setspace}
\setstretch{1.125}

%% URL styling - use normal text instead of monospace
\urlstyle{same}

%% Patch \httplink and \httpslink to handle full URLs
% The original moderncv macros always prepend http:// or https:// to URLs.
% This causes issues when the URL already has a protocol (e.g., https://example.com)
% resulting in malformed URLs like https://https://example.com.
% This patch checks if the URL already contains "://" and uses it directly if so.
% Note: We use \str_set:Nx (x-expansion) to expand macros like \@homepage first.
\ExplSyntaxOn
\str_new:N \g_yamlresume_url_str

\RenewDocumentCommand{\httpslink}{O{}m}{
  \str_set:Nx \g_yamlresume_url_str {#2}
  \str_if_in:NnTF \g_yamlresume_url_str {://}
    { \url{#2} }
    { \url{https://#2} }
}

\RenewDocumentCommand{\httplink}{O{}m}{
  \str_set:Nx \g_yamlresume_url_str {#2}
  \str_if_in:NnTF \g_yamlresume_url_str {://}
    { \url{#2} }
    { \url{http://#2} }
}
\ExplSyntaxOff

\name{Camille Moreau}{}
\title{Ingénieur mécanique senior spécialisé en conception énergétique}
\phone[mobile]{+33 6 12 34 56 78}
\email{camille.moreau@example.fr}
\homepage{https://www.camillemoreau-ingenierie.fr}

\address{24 rue des Ateliers, Lyon, Auvergne-Rhône-Alpes, France, 69003}{}{}

\extrainfo{{\small \faLinkedin}\ \href{https://www.linkedin.com/in/camille-moreau-meca}{@camille-moreau-meca} {} {} {} • {} {} {} 
{\small \faGithub}\ \href{https://github.com/camillemeca}{@camillemeca}}

\begin{document}

\maketitle

\section{Informations de base}

\cvline{}{\begin{itemize}
\item Ingénieur mécanique avec plus de 8 ans d'expérience dans la conception, la simulation et l'industrialisation de systèmes mécaniques et thermiques.
\item Expertise en CAO 3D, calcul de structures, CFD et gestion de projets industriels, du cahier des charges à la mise en production.
\item Solide expérience en environnement réglementé, coordination de fournisseurs et amélioration continue de la performance produit.
\item Reconnu pour la rigueur technique, la communication transversale et la capacité à transformer des besoins clients en solutions fiables.
\end{itemize}}
\section{Formation}

\cventry{sept. 2012–juil. 2017}
        {Master, Génie mécanique, Score : 3.7/4.0}
        {INSA Lyon}
        {\url{https://www.insa-lyon.fr}}
        {}
        {\begin{itemize}
\item Formation d'ingénieur généraliste orientée conception mécanique, énergétique et industrialisation.
\item Projets d'équipe sur des systèmes de transmission, des échangeurs thermiques et des bancs d'essai instrumentés.
\item Mémoire de fin d'études consacré à l'optimisation topologique de pièces mécaniques allégées.
\end{itemize}
\textbf{Cours} : Résistance des matériaux, Thermodynamique appliquée, Mécanique des fluides, Conception assistée par ordinateur, Fabrication mécanique, Vibrations et dynamique des structures, Science des matériaux, Automatique et systèmes embarqués}
\section{Expérience professionnelle}

\cventry{avr. 2021–Present}
        {Ingénieur mécanique senior}
        {Voltaire Énergie}
        {\url{https://www.voltaire-energie.fr}}
        {}
        {\begin{itemize}
\item Piloter la conception mécanique de modules de récupération de chaleur pour des sites industriels et des réseaux de chaleur urbains.
\item Coordonner les études thermiques, les calculs de résistance et la sélection des matériaux avec les équipes procédés et électriques.
\item Gérer les revues de conception, les analyses de risques et la documentation technique jusqu'à la réception client.
\item Encadrer deux ingénieurs juniors et animer les échanges avec les fournisseurs stratégiques.
\end{itemize}
\textbf{Mots-clés} : Conception mécanique, Thermique, CAO, Gestion de projet, Validation}

\cventry{sept. 2017–mars 2021}
        {Ingénieur mécanique}
        {Atelier Aérotech}
        {\url{https://www.atelier-aerotech.fr}}
        {}
        {\begin{itemize}
\item Concevoir des ensembles mécaniques et des supports de fixation pour des équipements aéronautiques.
\item Réaliser des simulations statiques et modales sous ANSYS afin de valider les choix de conception.
\item Collaborer avec l'atelier de fabrication pour optimiser les tolérances, les matériaux et les coûts de production.
\item Participer aux essais de qualification et rédiger les rapports techniques associés.
\end{itemize}
\textbf{Mots-clés} : Aéronautique, ANSYS, Tolérancement, Essais, Fabrication}

\cventry{févr. 2016–août 2016}
        {Ingénieur mécanique stagiaire}
        {Groupe Mécallia}
        {\url{https://www.groupe-mecallia.fr}}
        {}
        {\begin{itemize}
\item Contribuer à la refonte d'un banc d'essai de pompes industrielles.
\item Modéliser les pièces sous SolidWorks et produire les plans de fabrication.
\item Analyser les résultats d'essais et proposer des améliorations de montage.
\end{itemize}
\textbf{Mots-clés} : SolidWorks, Pompes, Banc d'essai, Analyse}
\section{Langues}

\cvline{Français}{Niveau natif ou bilingue \hfill \textbf{Mots-clés} : Langue maternelle}
\cvline{Anglais}{Niveau professionnel complet \hfill \textbf{Mots-clés} : TOEIC 920, Rédaction technique}
\cvline{Espagnol}{Niveau de travail limité \hfill \textbf{Mots-clés} : Conversation professionnelle}
\section{Compétences}

\cvline{Conception mécanique}{Expert \hfill \textbf{Mots-clés} : SolidWorks, CATIA V5, Tolérancement, RDM, DFM}
\cvline{Simulation numérique}{Avancé \hfill \textbf{Mots-clés} : ANSYS, CFD, Thermique, Fatigue, Matériaux composites}
\cvline{Gestion de projet}{Avancé \hfill \textbf{Mots-clés} : Planification, Budget, Fournisseurs, Agile, Documentation}
\cvline{Industrialisation}{Intermédiaire \hfill \textbf{Mots-clés} : Lean, Prototypage, Essais, Qualité, Amélioration continue}
\section{Récompenses}

\cventry{juin 2023}
        {Association Française de Mécanique}
        {Prix de l'innovation mécanique}
        {}
        {}
        {Distinction remise pour un module compact de récupération de chaleur réduisant les pertes thermiques de 18 \% sur une ligne de production.}
\section{Certificats}

\cventry{nov. 2019}
        {Dassault Systèmes}
        {Certification professionnelle SolidWorks}
        {\url{https://www.solidworks.com/certification}}
        {}
        {}

\cventry{mai 2022}
        {IASSC}
        {Ceinture verte Lean Six Sigma}
        {\url{https://www.iassc.org}}
        {}
        {}
\section{Publications}

\cventry{oct. 2022}
        {Optimisation thermique des échangeurs à plaques pour réseaux de chaleur}
        {Revue Française de Génie Mécanique}
        {\url{https://www.rfgm.fr/articles/echangeurs-plaques}}
        {}
        {\begin{itemize}
\item Article présentant une méthode de dimensionnement d'échangeurs à plaques intégrant les contraintes de maintenance.
\item Résultats issus d'une campagne de mesures sur trois sites industriels et d'un modèle CFD validé.
\end{itemize}}
\section{Projets}

\cventry{janv. 2022–sept. 2023}
        {Module compact de récupération de chaleur destiné aux procédés industriels.}
        {Récupérateur de chaleur industriel}
        {\url{https://www.voltaire-energie.fr/projets/recuperateur}}
        {}
        {\begin{itemize}
\item Conception mécanique complète d'un échangeur air-eau modulaire.
\item Intégration des contraintes de maintenance, d'isolation et de sécurité.
\item Validation par calculs thermomécaniques et essais sur banc.
\end{itemize}
\textbf{Mots-clés} : Thermique, Échangeur, Industrie, Modélisation}

\cventry{mars 2019–févr. 2020}
        {Banc instrumenté pour caractériser le rendement et la fiabilité de pompes industrielles.}
        {Banc d'essai pour pompes à haut rendement}
        {}
        {}
        {\begin{itemize}
\item Définition de l'architecture mécanique, des capteurs et de la chaîne d'acquisition.
\item Réalisation des plans, du dossier de fabrication et du protocole d'essai.
\item Réduction du temps de qualification des pompes de 25 \%.
\end{itemize}
\textbf{Mots-clés} : Banc d'essai, Instrumentation, Pompes, Essais}
\section{Centres d'intérêt}

\cvline{Sports}{Randonnée, Escalade, Cyclisme}
\cvline{Musique}{Guitare, Jazz}
\cvline{Fabrication}{Impression 3D, Menuiserie}
\section{Bénévolat}

\cventry{sept. 2018–Present}
        {Bénévole technique}
        {Ingénieurs Sans Frontières France}
        {\url{https://www.isf-france.org}}
        {}
        {\begin{itemize}
\item Apporter un appui technique à des projets d'accès à l'eau et d'énergie rurale.
\item Participer à la revue de plans, au dimensionnement d'équipements et à la formation des équipes locales.
\item Contribuer à la documentation et au retour d'expérience des missions.
\end{itemize}}
    
\end{document}